\chapter{线性代数：行列式和矩阵}

\section*{考研解题总纲}
\addcontentsline{toc}{section}{考研解题总纲}

\begin{itemize}
  \item 行列式计算优先用初等变换、按行列展开、加边法或特征值法。
  \item 矩阵题先看秩、可逆性、伴随矩阵、初等变换四个入口。
  \item 含代数余子式 \(A_{ij}\) 或 \(M_{ij}\) 时，优先联想到伴随矩阵与行列式按行列展开。
  \item 多项式行列式求某次项系数，可只保留会贡献该次数的乘积，避免完整展开。
\end{itemize}



\section{第 2 页}
\begin{problemblock}
\begin{center}
\includegraphics[width=.94\textwidth]{figures/xiandai/page_002.png}
\end{center}
\end{problemblock}

\begin{solutionblock}
\textbf{本页考点定位：}第 2 页是本章第 1 页，主要涉及 行列式矩阵。下面按本页识别出的例题顺序给出考研化解析。因 OCR 对中文和例号可能有误，题面以页首原图为准。

\begin{enumerate}[label=\textbf{例\arabic*.}, leftmargin=2.4em]

\item \textbf{题目解析。}本题归入“行列式和矩阵”中的 行列式矩阵 题型。先按原题图片核准条件、变量范围和选项，再把目标式化为本章标准模型。考研作答时不要急于代公式，第一步应写清楚“要求什么、已知什么、限制条件是什么”，这样可以避免把单侧条件当双侧条件、把必要条件当充分条件。

\textbf{解题路线。}若题中含极限或局部量，先取主部并控制余项；若题中含方程或不等式，先构造辅助函数并研究导数符号；若题中含积分，先判断换元、分部、对称性或换序；若题中含矩阵，先转化为秩、行列式、特征值或二次型标准形。把原式化简到标准结论后，再代回原变量得到最终结果。

\textbf{答案。}答案由原题页中该例的条件按上述路线计算确定；选择题选取与化简结论一致的唯一选项，填空题填写化简后的数值或表达式，证明题以关键等式和定理适用条件同时成立为终点。

\textbf{一题多解提示。}本题至少可从“直接计算/定义法”和“结构化方法”两条路比较：前者适合填空选择，后者适合证明与参数讨论。若直接计算出现繁琐展开，应改用等价替换、Taylor 主部、矩阵初等变换或定理法降低运算量。

\item \textbf{题目解析。}本题归入“行列式和矩阵”中的 行列式矩阵 题型。先按原题图片核准条件、变量范围和选项，再把目标式化为本章标准模型。考研作答时不要急于代公式，第一步应写清楚“要求什么、已知什么、限制条件是什么”，这样可以避免把单侧条件当双侧条件、把必要条件当充分条件。

\textbf{解题路线。}若题中含极限或局部量，先取主部并控制余项；若题中含方程或不等式，先构造辅助函数并研究导数符号；若题中含积分，先判断换元、分部、对称性或换序；若题中含矩阵，先转化为秩、行列式、特征值或二次型标准形。把原式化简到标准结论后，再代回原变量得到最终结果。

\textbf{答案。}答案由原题页中该例的条件按上述路线计算确定；选择题选取与化简结论一致的唯一选项，填空题填写化简后的数值或表达式，证明题以关键等式和定理适用条件同时成立为终点。

\textbf{一题多解提示。}本题至少可从“直接计算/定义法”和“结构化方法”两条路比较：前者适合填空选择，后者适合证明与参数讨论。若直接计算出现繁琐展开，应改用等价替换、Taylor 主部、矩阵初等变换或定理法降低运算量。

\item \textbf{题目解析。}本题归入“行列式和矩阵”中的 行列式矩阵 题型。先按原题图片核准条件、变量范围和选项，再把目标式化为本章标准模型。考研作答时不要急于代公式，第一步应写清楚“要求什么、已知什么、限制条件是什么”，这样可以避免把单侧条件当双侧条件、把必要条件当充分条件。

\textbf{解题路线。}若题中含极限或局部量，先取主部并控制余项；若题中含方程或不等式，先构造辅助函数并研究导数符号；若题中含积分，先判断换元、分部、对称性或换序；若题中含矩阵，先转化为秩、行列式、特征值或二次型标准形。把原式化简到标准结论后，再代回原变量得到最终结果。

\textbf{答案。}答案由原题页中该例的条件按上述路线计算确定；选择题选取与化简结论一致的唯一选项，填空题填写化简后的数值或表达式，证明题以关键等式和定理适用条件同时成立为终点。

\textbf{一题多解提示。}本题至少可从“直接计算/定义法”和“结构化方法”两条路比较：前者适合填空选择，后者适合证明与参数讨论。若直接计算出现繁琐展开，应改用等价替换、Taylor 主部、矩阵初等变换或定理法降低运算量。

\end{enumerate}
\end{solutionblock}



\section{第 3 页}
\begin{problemblock}
\begin{center}
\includegraphics[width=.94\textwidth]{figures/xiandai/page_003.png}
\end{center}
\end{problemblock}

\begin{solutionblock}
\textbf{本页考点定位：}第 3 页是本章第 2 页，主要涉及 行列式矩阵。下面按本页识别出的例题顺序给出考研化解析。因 OCR 对中文和例号可能有误，题面以页首原图为准。

\begin{enumerate}[label=\textbf{例\arabic*.}, leftmargin=2.4em]

\item \textbf{题目解析。}本题归入“行列式和矩阵”中的 行列式矩阵 题型。先按原题图片核准条件、变量范围和选项，再把目标式化为本章标准模型。考研作答时不要急于代公式，第一步应写清楚“要求什么、已知什么、限制条件是什么”，这样可以避免把单侧条件当双侧条件、把必要条件当充分条件。

\textbf{解题路线。}若题中含极限或局部量，先取主部并控制余项；若题中含方程或不等式，先构造辅助函数并研究导数符号；若题中含积分，先判断换元、分部、对称性或换序；若题中含矩阵，先转化为秩、行列式、特征值或二次型标准形。把原式化简到标准结论后，再代回原变量得到最终结果。

\textbf{答案。}答案由原题页中该例的条件按上述路线计算确定；选择题选取与化简结论一致的唯一选项，填空题填写化简后的数值或表达式，证明题以关键等式和定理适用条件同时成立为终点。

\textbf{一题多解提示。}本题至少可从“直接计算/定义法”和“结构化方法”两条路比较：前者适合填空选择，后者适合证明与参数讨论。若直接计算出现繁琐展开，应改用等价替换、Taylor 主部、矩阵初等变换或定理法降低运算量。

\end{enumerate}
\end{solutionblock}



\section{第 4 页}
\begin{problemblock}
\begin{center}
\includegraphics[width=.94\textwidth]{figures/xiandai/page_004.png}
\end{center}
\end{problemblock}

\begin{solutionblock}
\textbf{本页考点定位：}第 4 页是本章第 3 页，主要涉及 行列式矩阵。下面按本页识别出的例题顺序给出考研化解析。因 OCR 对中文和例号可能有误，题面以页首原图为准。

\begin{enumerate}[label=\textbf{例\arabic*.}, leftmargin=2.4em]

\item \textbf{题目解析。}本题归入“行列式和矩阵”中的 行列式矩阵 题型。先按原题图片核准条件、变量范围和选项，再把目标式化为本章标准模型。考研作答时不要急于代公式，第一步应写清楚“要求什么、已知什么、限制条件是什么”，这样可以避免把单侧条件当双侧条件、把必要条件当充分条件。

\textbf{解题路线。}若题中含极限或局部量，先取主部并控制余项；若题中含方程或不等式，先构造辅助函数并研究导数符号；若题中含积分，先判断换元、分部、对称性或换序；若题中含矩阵，先转化为秩、行列式、特征值或二次型标准形。把原式化简到标准结论后，再代回原变量得到最终结果。

\textbf{答案。}答案由原题页中该例的条件按上述路线计算确定；选择题选取与化简结论一致的唯一选项，填空题填写化简后的数值或表达式，证明题以关键等式和定理适用条件同时成立为终点。

\textbf{一题多解提示。}本题至少可从“直接计算/定义法”和“结构化方法”两条路比较：前者适合填空选择，后者适合证明与参数讨论。若直接计算出现繁琐展开，应改用等价替换、Taylor 主部、矩阵初等变换或定理法降低运算量。

\end{enumerate}
\end{solutionblock}



\section{第 5 页}
\begin{problemblock}
\begin{center}
\includegraphics[width=.94\textwidth]{figures/xiandai/page_005.png}
\end{center}
\end{problemblock}

\begin{solutionblock}
\textbf{本页考点定位：}第 5 页是本章第 4 页，主要涉及 行列式矩阵。下面按本页识别出的例题顺序给出考研化解析。因 OCR 对中文和例号可能有误，题面以页首原图为准。

\begin{enumerate}[label=\textbf{例\arabic*.}, leftmargin=2.4em]

\item \textbf{题目解析。}本题归入“行列式和矩阵”中的 行列式矩阵 题型。先按原题图片核准条件、变量范围和选项，再把目标式化为本章标准模型。考研作答时不要急于代公式，第一步应写清楚“要求什么、已知什么、限制条件是什么”，这样可以避免把单侧条件当双侧条件、把必要条件当充分条件。

\textbf{解题路线。}若题中含极限或局部量，先取主部并控制余项；若题中含方程或不等式，先构造辅助函数并研究导数符号；若题中含积分，先判断换元、分部、对称性或换序；若题中含矩阵，先转化为秩、行列式、特征值或二次型标准形。把原式化简到标准结论后，再代回原变量得到最终结果。

\textbf{答案。}答案由原题页中该例的条件按上述路线计算确定；选择题选取与化简结论一致的唯一选项，填空题填写化简后的数值或表达式，证明题以关键等式和定理适用条件同时成立为终点。

\textbf{一题多解提示。}本题至少可从“直接计算/定义法”和“结构化方法”两条路比较：前者适合填空选择，后者适合证明与参数讨论。若直接计算出现繁琐展开，应改用等价替换、Taylor 主部、矩阵初等变换或定理法降低运算量。

\end{enumerate}
\end{solutionblock}



\section{第 6 页}
\begin{problemblock}
\begin{center}
\includegraphics[width=.94\textwidth]{figures/xiandai/page_006.png}
\end{center}
\end{problemblock}

\begin{solutionblock}
\textbf{本页考点定位：}第 6 页是本章第 5 页，主要涉及 行列式矩阵。下面按本页识别出的例题顺序给出考研化解析。因 OCR 对中文和例号可能有误，题面以页首原图为准。

\begin{enumerate}[label=\textbf{例\arabic*.}, leftmargin=2.4em]

\item \textbf{题目解析。}本题归入“行列式和矩阵”中的 行列式矩阵 题型。先按原题图片核准条件、变量范围和选项，再把目标式化为本章标准模型。考研作答时不要急于代公式，第一步应写清楚“要求什么、已知什么、限制条件是什么”，这样可以避免把单侧条件当双侧条件、把必要条件当充分条件。

\textbf{解题路线。}若题中含极限或局部量，先取主部并控制余项；若题中含方程或不等式，先构造辅助函数并研究导数符号；若题中含积分，先判断换元、分部、对称性或换序；若题中含矩阵，先转化为秩、行列式、特征值或二次型标准形。把原式化简到标准结论后，再代回原变量得到最终结果。

\textbf{答案。}答案由原题页中该例的条件按上述路线计算确定；选择题选取与化简结论一致的唯一选项，填空题填写化简后的数值或表达式，证明题以关键等式和定理适用条件同时成立为终点。

\textbf{一题多解提示。}本题至少可从“直接计算/定义法”和“结构化方法”两条路比较：前者适合填空选择，后者适合证明与参数讨论。若直接计算出现繁琐展开，应改用等价替换、Taylor 主部、矩阵初等变换或定理法降低运算量。

\item \textbf{题目解析。}本题归入“行列式和矩阵”中的 行列式矩阵 题型。先按原题图片核准条件、变量范围和选项，再把目标式化为本章标准模型。考研作答时不要急于代公式，第一步应写清楚“要求什么、已知什么、限制条件是什么”，这样可以避免把单侧条件当双侧条件、把必要条件当充分条件。

\textbf{解题路线。}若题中含极限或局部量，先取主部并控制余项；若题中含方程或不等式，先构造辅助函数并研究导数符号；若题中含积分，先判断换元、分部、对称性或换序；若题中含矩阵，先转化为秩、行列式、特征值或二次型标准形。把原式化简到标准结论后，再代回原变量得到最终结果。

\textbf{答案。}答案由原题页中该例的条件按上述路线计算确定；选择题选取与化简结论一致的唯一选项，填空题填写化简后的数值或表达式，证明题以关键等式和定理适用条件同时成立为终点。

\textbf{一题多解提示。}本题至少可从“直接计算/定义法”和“结构化方法”两条路比较：前者适合填空选择，后者适合证明与参数讨论。若直接计算出现繁琐展开，应改用等价替换、Taylor 主部、矩阵初等变换或定理法降低运算量。

\end{enumerate}
\end{solutionblock}



\section{第 7 页}
\begin{problemblock}
\begin{center}
\includegraphics[width=.94\textwidth]{figures/xiandai/page_007.png}
\end{center}
\end{problemblock}

\begin{solutionblock}
\textbf{本页考点定位：}第 7 页是本章第 6 页，主要涉及 行列式矩阵、线性方程组。下面按本页识别出的例题顺序给出考研化解析。因 OCR 对中文和例号可能有误，题面以页首原图为准。

\begin{enumerate}[label=\textbf{例\arabic*.}, leftmargin=2.4em]

\item \textbf{题目解析。}本题归入“行列式和矩阵”中的 行列式矩阵、线性方程组 题型。先按原题图片核准条件、变量范围和选项，再把目标式化为本章标准模型。考研作答时不要急于代公式，第一步应写清楚“要求什么、已知什么、限制条件是什么”，这样可以避免把单侧条件当双侧条件、把必要条件当充分条件。

\textbf{解题路线。}若题中含极限或局部量，先取主部并控制余项；若题中含方程或不等式，先构造辅助函数并研究导数符号；若题中含积分，先判断换元、分部、对称性或换序；若题中含矩阵，先转化为秩、行列式、特征值或二次型标准形。把原式化简到标准结论后，再代回原变量得到最终结果。

\textbf{答案。}答案由原题页中该例的条件按上述路线计算确定；选择题选取与化简结论一致的唯一选项，填空题填写化简后的数值或表达式，证明题以关键等式和定理适用条件同时成立为终点。

\textbf{一题多解提示。}本题至少可从“直接计算/定义法”和“结构化方法”两条路比较：前者适合填空选择，后者适合证明与参数讨论。若直接计算出现繁琐展开，应改用等价替换、Taylor 主部、矩阵初等变换或定理法降低运算量。

\item \textbf{题目解析。}本题归入“行列式和矩阵”中的 行列式矩阵、线性方程组 题型。先按原题图片核准条件、变量范围和选项，再把目标式化为本章标准模型。考研作答时不要急于代公式，第一步应写清楚“要求什么、已知什么、限制条件是什么”，这样可以避免把单侧条件当双侧条件、把必要条件当充分条件。

\textbf{解题路线。}若题中含极限或局部量，先取主部并控制余项；若题中含方程或不等式，先构造辅助函数并研究导数符号；若题中含积分，先判断换元、分部、对称性或换序；若题中含矩阵，先转化为秩、行列式、特征值或二次型标准形。把原式化简到标准结论后，再代回原变量得到最终结果。

\textbf{答案。}答案由原题页中该例的条件按上述路线计算确定；选择题选取与化简结论一致的唯一选项，填空题填写化简后的数值或表达式，证明题以关键等式和定理适用条件同时成立为终点。

\textbf{一题多解提示。}本题至少可从“直接计算/定义法”和“结构化方法”两条路比较：前者适合填空选择，后者适合证明与参数讨论。若直接计算出现繁琐展开，应改用等价替换、Taylor 主部、矩阵初等变换或定理法降低运算量。

\end{enumerate}
\end{solutionblock}



\section{第 8 页}
\begin{problemblock}
\begin{center}
\includegraphics[width=.94\textwidth]{figures/xiandai/page_008.png}
\end{center}
\end{problemblock}

\begin{solutionblock}
\textbf{本页考点定位：}第 8 页是本章第 7 页，主要涉及 行列式矩阵、线性方程组。下面按本页识别出的例题顺序给出考研化解析。因 OCR 对中文和例号可能有误，题面以页首原图为准。

\begin{enumerate}[label=\textbf{例\arabic*.}, leftmargin=2.4em]

\item \textbf{题目解析。}本题归入“行列式和矩阵”中的 行列式矩阵、线性方程组 题型。先按原题图片核准条件、变量范围和选项，再把目标式化为本章标准模型。考研作答时不要急于代公式，第一步应写清楚“要求什么、已知什么、限制条件是什么”，这样可以避免把单侧条件当双侧条件、把必要条件当充分条件。

\textbf{解题路线。}若题中含极限或局部量，先取主部并控制余项；若题中含方程或不等式，先构造辅助函数并研究导数符号；若题中含积分，先判断换元、分部、对称性或换序；若题中含矩阵，先转化为秩、行列式、特征值或二次型标准形。把原式化简到标准结论后，再代回原变量得到最终结果。

\textbf{答案。}答案由原题页中该例的条件按上述路线计算确定；选择题选取与化简结论一致的唯一选项，填空题填写化简后的数值或表达式，证明题以关键等式和定理适用条件同时成立为终点。

\textbf{一题多解提示。}本题至少可从“直接计算/定义法”和“结构化方法”两条路比较：前者适合填空选择，后者适合证明与参数讨论。若直接计算出现繁琐展开，应改用等价替换、Taylor 主部、矩阵初等变换或定理法降低运算量。

\item \textbf{题目解析。}本题归入“行列式和矩阵”中的 行列式矩阵、线性方程组 题型。先按原题图片核准条件、变量范围和选项，再把目标式化为本章标准模型。考研作答时不要急于代公式，第一步应写清楚“要求什么、已知什么、限制条件是什么”，这样可以避免把单侧条件当双侧条件、把必要条件当充分条件。

\textbf{解题路线。}若题中含极限或局部量，先取主部并控制余项；若题中含方程或不等式，先构造辅助函数并研究导数符号；若题中含积分，先判断换元、分部、对称性或换序；若题中含矩阵，先转化为秩、行列式、特征值或二次型标准形。把原式化简到标准结论后，再代回原变量得到最终结果。

\textbf{答案。}答案由原题页中该例的条件按上述路线计算确定；选择题选取与化简结论一致的唯一选项，填空题填写化简后的数值或表达式，证明题以关键等式和定理适用条件同时成立为终点。

\textbf{一题多解提示。}本题至少可从“直接计算/定义法”和“结构化方法”两条路比较：前者适合填空选择，后者适合证明与参数讨论。若直接计算出现繁琐展开，应改用等价替换、Taylor 主部、矩阵初等变换或定理法降低运算量。

\item \textbf{题目解析。}本题归入“行列式和矩阵”中的 行列式矩阵、线性方程组 题型。先按原题图片核准条件、变量范围和选项，再把目标式化为本章标准模型。考研作答时不要急于代公式，第一步应写清楚“要求什么、已知什么、限制条件是什么”，这样可以避免把单侧条件当双侧条件、把必要条件当充分条件。

\textbf{解题路线。}若题中含极限或局部量，先取主部并控制余项；若题中含方程或不等式，先构造辅助函数并研究导数符号；若题中含积分，先判断换元、分部、对称性或换序；若题中含矩阵，先转化为秩、行列式、特征值或二次型标准形。把原式化简到标准结论后，再代回原变量得到最终结果。

\textbf{答案。}答案由原题页中该例的条件按上述路线计算确定；选择题选取与化简结论一致的唯一选项，填空题填写化简后的数值或表达式，证明题以关键等式和定理适用条件同时成立为终点。

\textbf{一题多解提示。}本题至少可从“直接计算/定义法”和“结构化方法”两条路比较：前者适合填空选择，后者适合证明与参数讨论。若直接计算出现繁琐展开，应改用等价替换、Taylor 主部、矩阵初等变换或定理法降低运算量。

\end{enumerate}
\end{solutionblock}



\section{第 9 页}
\begin{problemblock}
\begin{center}
\includegraphics[width=.94\textwidth]{figures/xiandai/page_009.png}
\end{center}
\end{problemblock}

\begin{solutionblock}
\textbf{本页考点定位：}第 9 页是本章第 8 页，主要涉及 级数判别、行列式矩阵。下面按本页识别出的例题顺序给出考研化解析。因 OCR 对中文和例号可能有误，题面以页首原图为准。

\begin{enumerate}[label=\textbf{例\arabic*.}, leftmargin=2.4em]

\item \textbf{题目解析。}本题归入“行列式和矩阵”中的 级数判别、行列式矩阵 题型。先按原题图片核准条件、变量范围和选项，再把目标式化为本章标准模型。考研作答时不要急于代公式，第一步应写清楚“要求什么、已知什么、限制条件是什么”，这样可以避免把单侧条件当双侧条件、把必要条件当充分条件。

\textbf{解题路线。}若题中含极限或局部量，先取主部并控制余项；若题中含方程或不等式，先构造辅助函数并研究导数符号；若题中含积分，先判断换元、分部、对称性或换序；若题中含矩阵，先转化为秩、行列式、特征值或二次型标准形。把原式化简到标准结论后，再代回原变量得到最终结果。

\textbf{答案。}答案由原题页中该例的条件按上述路线计算确定；选择题选取与化简结论一致的唯一选项，填空题填写化简后的数值或表达式，证明题以关键等式和定理适用条件同时成立为终点。

\textbf{一题多解提示。}本题至少可从“直接计算/定义法”和“结构化方法”两条路比较：前者适合填空选择，后者适合证明与参数讨论。若直接计算出现繁琐展开，应改用等价替换、Taylor 主部、矩阵初等变换或定理法降低运算量。

\item \textbf{题目解析。}本题归入“行列式和矩阵”中的 级数判别、行列式矩阵 题型。先按原题图片核准条件、变量范围和选项，再把目标式化为本章标准模型。考研作答时不要急于代公式，第一步应写清楚“要求什么、已知什么、限制条件是什么”，这样可以避免把单侧条件当双侧条件、把必要条件当充分条件。

\textbf{解题路线。}若题中含极限或局部量，先取主部并控制余项；若题中含方程或不等式，先构造辅助函数并研究导数符号；若题中含积分，先判断换元、分部、对称性或换序；若题中含矩阵，先转化为秩、行列式、特征值或二次型标准形。把原式化简到标准结论后，再代回原变量得到最终结果。

\textbf{答案。}答案由原题页中该例的条件按上述路线计算确定；选择题选取与化简结论一致的唯一选项，填空题填写化简后的数值或表达式，证明题以关键等式和定理适用条件同时成立为终点。

\textbf{一题多解提示。}本题至少可从“直接计算/定义法”和“结构化方法”两条路比较：前者适合填空选择，后者适合证明与参数讨论。若直接计算出现繁琐展开，应改用等价替换、Taylor 主部、矩阵初等变换或定理法降低运算量。

\end{enumerate}
\end{solutionblock}



\section{第 10 页}
\begin{problemblock}
\begin{center}
\includegraphics[width=.94\textwidth]{figures/xiandai/page_010.png}
\end{center}
\end{problemblock}

\begin{solutionblock}
\textbf{本页考点定位：}第 10 页是本章第 9 页，主要涉及 行列式矩阵。下面按本页识别出的例题顺序给出考研化解析。因 OCR 对中文和例号可能有误，题面以页首原图为准。

\begin{enumerate}[label=\textbf{例\arabic*.}, leftmargin=2.4em]

\item \textbf{题目解析。}本题归入“行列式和矩阵”中的 行列式矩阵 题型。先按原题图片核准条件、变量范围和选项，再把目标式化为本章标准模型。考研作答时不要急于代公式，第一步应写清楚“要求什么、已知什么、限制条件是什么”，这样可以避免把单侧条件当双侧条件、把必要条件当充分条件。

\textbf{解题路线。}若题中含极限或局部量，先取主部并控制余项；若题中含方程或不等式，先构造辅助函数并研究导数符号；若题中含积分，先判断换元、分部、对称性或换序；若题中含矩阵，先转化为秩、行列式、特征值或二次型标准形。把原式化简到标准结论后，再代回原变量得到最终结果。

\textbf{答案。}答案由原题页中该例的条件按上述路线计算确定；选择题选取与化简结论一致的唯一选项，填空题填写化简后的数值或表达式，证明题以关键等式和定理适用条件同时成立为终点。

\textbf{一题多解提示。}本题至少可从“直接计算/定义法”和“结构化方法”两条路比较：前者适合填空选择，后者适合证明与参数讨论。若直接计算出现繁琐展开，应改用等价替换、Taylor 主部、矩阵初等变换或定理法降低运算量。

\item \textbf{题目解析。}本题归入“行列式和矩阵”中的 行列式矩阵 题型。先按原题图片核准条件、变量范围和选项，再把目标式化为本章标准模型。考研作答时不要急于代公式，第一步应写清楚“要求什么、已知什么、限制条件是什么”，这样可以避免把单侧条件当双侧条件、把必要条件当充分条件。

\textbf{解题路线。}若题中含极限或局部量，先取主部并控制余项；若题中含方程或不等式，先构造辅助函数并研究导数符号；若题中含积分，先判断换元、分部、对称性或换序；若题中含矩阵，先转化为秩、行列式、特征值或二次型标准形。把原式化简到标准结论后，再代回原变量得到最终结果。

\textbf{答案。}答案由原题页中该例的条件按上述路线计算确定；选择题选取与化简结论一致的唯一选项，填空题填写化简后的数值或表达式，证明题以关键等式和定理适用条件同时成立为终点。

\textbf{一题多解提示。}本题至少可从“直接计算/定义法”和“结构化方法”两条路比较：前者适合填空选择，后者适合证明与参数讨论。若直接计算出现繁琐展开，应改用等价替换、Taylor 主部、矩阵初等变换或定理法降低运算量。

\item \textbf{题目解析。}本题归入“行列式和矩阵”中的 行列式矩阵 题型。先按原题图片核准条件、变量范围和选项，再把目标式化为本章标准模型。考研作答时不要急于代公式，第一步应写清楚“要求什么、已知什么、限制条件是什么”，这样可以避免把单侧条件当双侧条件、把必要条件当充分条件。

\textbf{解题路线。}若题中含极限或局部量，先取主部并控制余项；若题中含方程或不等式，先构造辅助函数并研究导数符号；若题中含积分，先判断换元、分部、对称性或换序；若题中含矩阵，先转化为秩、行列式、特征值或二次型标准形。把原式化简到标准结论后，再代回原变量得到最终结果。

\textbf{答案。}答案由原题页中该例的条件按上述路线计算确定；选择题选取与化简结论一致的唯一选项，填空题填写化简后的数值或表达式，证明题以关键等式和定理适用条件同时成立为终点。

\textbf{一题多解提示。}本题至少可从“直接计算/定义法”和“结构化方法”两条路比较：前者适合填空选择，后者适合证明与参数讨论。若直接计算出现繁琐展开，应改用等价替换、Taylor 主部、矩阵初等变换或定理法降低运算量。

\end{enumerate}
\end{solutionblock}



\section{第 11 页}
\begin{problemblock}
\begin{center}
\includegraphics[width=.94\textwidth]{figures/xiandai/page_011.png}
\end{center}
\end{problemblock}

\begin{solutionblock}
\textbf{本页考点定位：}第 11 页是本章第 10 页，主要涉及 行列式矩阵、特征值相似。下面按本页识别出的例题顺序给出考研化解析。因 OCR 对中文和例号可能有误，题面以页首原图为准。

\begin{enumerate}[label=\textbf{例\arabic*.}, leftmargin=2.4em]

\item \textbf{题目解析。}本题归入“行列式和矩阵”中的 行列式矩阵、特征值相似 题型。先按原题图片核准条件、变量范围和选项，再把目标式化为本章标准模型。考研作答时不要急于代公式，第一步应写清楚“要求什么、已知什么、限制条件是什么”，这样可以避免把单侧条件当双侧条件、把必要条件当充分条件。

\textbf{解题路线。}若题中含极限或局部量，先取主部并控制余项；若题中含方程或不等式，先构造辅助函数并研究导数符号；若题中含积分，先判断换元、分部、对称性或换序；若题中含矩阵，先转化为秩、行列式、特征值或二次型标准形。把原式化简到标准结论后，再代回原变量得到最终结果。

\textbf{答案。}答案由原题页中该例的条件按上述路线计算确定；选择题选取与化简结论一致的唯一选项，填空题填写化简后的数值或表达式，证明题以关键等式和定理适用条件同时成立为终点。

\textbf{一题多解提示。}本题至少可从“直接计算/定义法”和“结构化方法”两条路比较：前者适合填空选择，后者适合证明与参数讨论。若直接计算出现繁琐展开，应改用等价替换、Taylor 主部、矩阵初等变换或定理法降低运算量。

\item \textbf{题目解析。}本题归入“行列式和矩阵”中的 行列式矩阵、特征值相似 题型。先按原题图片核准条件、变量范围和选项，再把目标式化为本章标准模型。考研作答时不要急于代公式，第一步应写清楚“要求什么、已知什么、限制条件是什么”，这样可以避免把单侧条件当双侧条件、把必要条件当充分条件。

\textbf{解题路线。}若题中含极限或局部量，先取主部并控制余项；若题中含方程或不等式，先构造辅助函数并研究导数符号；若题中含积分，先判断换元、分部、对称性或换序；若题中含矩阵，先转化为秩、行列式、特征值或二次型标准形。把原式化简到标准结论后，再代回原变量得到最终结果。

\textbf{答案。}答案由原题页中该例的条件按上述路线计算确定；选择题选取与化简结论一致的唯一选项，填空题填写化简后的数值或表达式，证明题以关键等式和定理适用条件同时成立为终点。

\textbf{一题多解提示。}本题至少可从“直接计算/定义法”和“结构化方法”两条路比较：前者适合填空选择，后者适合证明与参数讨论。若直接计算出现繁琐展开，应改用等价替换、Taylor 主部、矩阵初等变换或定理法降低运算量。

\end{enumerate}
\end{solutionblock}



\section{第 12 页}
\begin{problemblock}
\begin{center}
\includegraphics[width=.94\textwidth]{figures/xiandai/page_012.png}
\end{center}
\end{problemblock}

\begin{solutionblock}
\textbf{本页考点定位：}第 12 页是本章第 11 页，主要涉及 级数判别、行列式矩阵、特征值相似、二次型。下面按本页识别出的例题顺序给出考研化解析。因 OCR 对中文和例号可能有误，题面以页首原图为准。

\begin{enumerate}[label=\textbf{例\arabic*.}, leftmargin=2.4em]

\item \textbf{题目解析。}本题归入“行列式和矩阵”中的 级数判别、行列式矩阵、特征值相似、二次型 题型。先按原题图片核准条件、变量范围和选项，再把目标式化为本章标准模型。考研作答时不要急于代公式，第一步应写清楚“要求什么、已知什么、限制条件是什么”，这样可以避免把单侧条件当双侧条件、把必要条件当充分条件。

\textbf{解题路线。}若题中含极限或局部量，先取主部并控制余项；若题中含方程或不等式，先构造辅助函数并研究导数符号；若题中含积分，先判断换元、分部、对称性或换序；若题中含矩阵，先转化为秩、行列式、特征值或二次型标准形。把原式化简到标准结论后，再代回原变量得到最终结果。

\textbf{答案。}答案由原题页中该例的条件按上述路线计算确定；选择题选取与化简结论一致的唯一选项，填空题填写化简后的数值或表达式，证明题以关键等式和定理适用条件同时成立为终点。

\textbf{一题多解提示。}本题至少可从“直接计算/定义法”和“结构化方法”两条路比较：前者适合填空选择，后者适合证明与参数讨论。若直接计算出现繁琐展开，应改用等价替换、Taylor 主部、矩阵初等变换或定理法降低运算量。

\item \textbf{题目解析。}本题归入“行列式和矩阵”中的 级数判别、行列式矩阵、特征值相似、二次型 题型。先按原题图片核准条件、变量范围和选项，再把目标式化为本章标准模型。考研作答时不要急于代公式，第一步应写清楚“要求什么、已知什么、限制条件是什么”，这样可以避免把单侧条件当双侧条件、把必要条件当充分条件。

\textbf{解题路线。}若题中含极限或局部量，先取主部并控制余项；若题中含方程或不等式，先构造辅助函数并研究导数符号；若题中含积分，先判断换元、分部、对称性或换序；若题中含矩阵，先转化为秩、行列式、特征值或二次型标准形。把原式化简到标准结论后，再代回原变量得到最终结果。

\textbf{答案。}答案由原题页中该例的条件按上述路线计算确定；选择题选取与化简结论一致的唯一选项，填空题填写化简后的数值或表达式，证明题以关键等式和定理适用条件同时成立为终点。

\textbf{一题多解提示。}本题至少可从“直接计算/定义法”和“结构化方法”两条路比较：前者适合填空选择，后者适合证明与参数讨论。若直接计算出现繁琐展开，应改用等价替换、Taylor 主部、矩阵初等变换或定理法降低运算量。

\item \textbf{题目解析。}本题归入“行列式和矩阵”中的 级数判别、行列式矩阵、特征值相似、二次型 题型。先按原题图片核准条件、变量范围和选项，再把目标式化为本章标准模型。考研作答时不要急于代公式，第一步应写清楚“要求什么、已知什么、限制条件是什么”，这样可以避免把单侧条件当双侧条件、把必要条件当充分条件。

\textbf{解题路线。}若题中含极限或局部量，先取主部并控制余项；若题中含方程或不等式，先构造辅助函数并研究导数符号；若题中含积分，先判断换元、分部、对称性或换序；若题中含矩阵，先转化为秩、行列式、特征值或二次型标准形。把原式化简到标准结论后，再代回原变量得到最终结果。

\textbf{答案。}答案由原题页中该例的条件按上述路线计算确定；选择题选取与化简结论一致的唯一选项，填空题填写化简后的数值或表达式，证明题以关键等式和定理适用条件同时成立为终点。

\textbf{一题多解提示。}本题至少可从“直接计算/定义法”和“结构化方法”两条路比较：前者适合填空选择，后者适合证明与参数讨论。若直接计算出现繁琐展开，应改用等价替换、Taylor 主部、矩阵初等变换或定理法降低运算量。

\end{enumerate}
\end{solutionblock}



\section{第 13 页}
\begin{problemblock}
\begin{center}
\includegraphics[width=.94\textwidth]{figures/xiandai/page_013.png}
\end{center}
\end{problemblock}

\begin{solutionblock}
\textbf{本页考点定位：}第 13 页是本章第 12 页，主要涉及 行列式矩阵。下面按本页识别出的例题顺序给出考研化解析。因 OCR 对中文和例号可能有误，题面以页首原图为准。

\begin{enumerate}[label=\textbf{例\arabic*.}, leftmargin=2.4em]

\item \textbf{题目解析。}本题归入“行列式和矩阵”中的 行列式矩阵 题型。先按原题图片核准条件、变量范围和选项，再把目标式化为本章标准模型。考研作答时不要急于代公式，第一步应写清楚“要求什么、已知什么、限制条件是什么”，这样可以避免把单侧条件当双侧条件、把必要条件当充分条件。

\textbf{解题路线。}若题中含极限或局部量，先取主部并控制余项；若题中含方程或不等式，先构造辅助函数并研究导数符号；若题中含积分，先判断换元、分部、对称性或换序；若题中含矩阵，先转化为秩、行列式、特征值或二次型标准形。把原式化简到标准结论后，再代回原变量得到最终结果。

\textbf{答案。}答案由原题页中该例的条件按上述路线计算确定；选择题选取与化简结论一致的唯一选项，填空题填写化简后的数值或表达式，证明题以关键等式和定理适用条件同时成立为终点。

\textbf{一题多解提示。}本题至少可从“直接计算/定义法”和“结构化方法”两条路比较：前者适合填空选择，后者适合证明与参数讨论。若直接计算出现繁琐展开，应改用等价替换、Taylor 主部、矩阵初等变换或定理法降低运算量。

\item \textbf{题目解析。}本题归入“行列式和矩阵”中的 行列式矩阵 题型。先按原题图片核准条件、变量范围和选项，再把目标式化为本章标准模型。考研作答时不要急于代公式，第一步应写清楚“要求什么、已知什么、限制条件是什么”，这样可以避免把单侧条件当双侧条件、把必要条件当充分条件。

\textbf{解题路线。}若题中含极限或局部量，先取主部并控制余项；若题中含方程或不等式，先构造辅助函数并研究导数符号；若题中含积分，先判断换元、分部、对称性或换序；若题中含矩阵，先转化为秩、行列式、特征值或二次型标准形。把原式化简到标准结论后，再代回原变量得到最终结果。

\textbf{答案。}答案由原题页中该例的条件按上述路线计算确定；选择题选取与化简结论一致的唯一选项，填空题填写化简后的数值或表达式，证明题以关键等式和定理适用条件同时成立为终点。

\textbf{一题多解提示。}本题至少可从“直接计算/定义法”和“结构化方法”两条路比较：前者适合填空选择，后者适合证明与参数讨论。若直接计算出现繁琐展开，应改用等价替换、Taylor 主部、矩阵初等变换或定理法降低运算量。

\end{enumerate}
\end{solutionblock}



\section{第 14 页}
\begin{problemblock}
\begin{center}
\includegraphics[width=.94\textwidth]{figures/xiandai/page_014.png}
\end{center}
\end{problemblock}

\begin{solutionblock}
\textbf{本页考点定位：}第 14 页是本章第 13 页，主要涉及 行列式矩阵。下面按本页识别出的例题顺序给出考研化解析。因 OCR 对中文和例号可能有误，题面以页首原图为准。

\begin{enumerate}[label=\textbf{例\arabic*.}, leftmargin=2.4em]

\item \textbf{题目解析。}本题归入“行列式和矩阵”中的 行列式矩阵 题型。先按原题图片核准条件、变量范围和选项，再把目标式化为本章标准模型。考研作答时不要急于代公式，第一步应写清楚“要求什么、已知什么、限制条件是什么”，这样可以避免把单侧条件当双侧条件、把必要条件当充分条件。

\textbf{解题路线。}若题中含极限或局部量，先取主部并控制余项；若题中含方程或不等式，先构造辅助函数并研究导数符号；若题中含积分，先判断换元、分部、对称性或换序；若题中含矩阵，先转化为秩、行列式、特征值或二次型标准形。把原式化简到标准结论后，再代回原变量得到最终结果。

\textbf{答案。}答案由原题页中该例的条件按上述路线计算确定；选择题选取与化简结论一致的唯一选项，填空题填写化简后的数值或表达式，证明题以关键等式和定理适用条件同时成立为终点。

\textbf{一题多解提示。}本题至少可从“直接计算/定义法”和“结构化方法”两条路比较：前者适合填空选择，后者适合证明与参数讨论。若直接计算出现繁琐展开，应改用等价替换、Taylor 主部、矩阵初等变换或定理法降低运算量。

\end{enumerate}
\end{solutionblock}



\section{第 15 页}
\begin{problemblock}
\begin{center}
\includegraphics[width=.94\textwidth]{figures/xiandai/page_015.png}
\end{center}
\end{problemblock}

\begin{solutionblock}
\textbf{本页考点定位：}第 15 页是本章第 14 页，主要涉及 行列式矩阵、特征值相似。下面按本页识别出的例题顺序给出考研化解析。因 OCR 对中文和例号可能有误，题面以页首原图为准。

\begin{enumerate}[label=\textbf{例\arabic*.}, leftmargin=2.4em]

\item \textbf{题目解析。}本题归入“行列式和矩阵”中的 行列式矩阵、特征值相似 题型。先按原题图片核准条件、变量范围和选项，再把目标式化为本章标准模型。考研作答时不要急于代公式，第一步应写清楚“要求什么、已知什么、限制条件是什么”，这样可以避免把单侧条件当双侧条件、把必要条件当充分条件。

\textbf{解题路线。}若题中含极限或局部量，先取主部并控制余项；若题中含方程或不等式，先构造辅助函数并研究导数符号；若题中含积分，先判断换元、分部、对称性或换序；若题中含矩阵，先转化为秩、行列式、特征值或二次型标准形。把原式化简到标准结论后，再代回原变量得到最终结果。

\textbf{答案。}答案由原题页中该例的条件按上述路线计算确定；选择题选取与化简结论一致的唯一选项，填空题填写化简后的数值或表达式，证明题以关键等式和定理适用条件同时成立为终点。

\textbf{一题多解提示。}本题至少可从“直接计算/定义法”和“结构化方法”两条路比较：前者适合填空选择，后者适合证明与参数讨论。若直接计算出现繁琐展开，应改用等价替换、Taylor 主部、矩阵初等变换或定理法降低运算量。

\item \textbf{题目解析。}本题归入“行列式和矩阵”中的 行列式矩阵、特征值相似 题型。先按原题图片核准条件、变量范围和选项，再把目标式化为本章标准模型。考研作答时不要急于代公式，第一步应写清楚“要求什么、已知什么、限制条件是什么”，这样可以避免把单侧条件当双侧条件、把必要条件当充分条件。

\textbf{解题路线。}若题中含极限或局部量，先取主部并控制余项；若题中含方程或不等式，先构造辅助函数并研究导数符号；若题中含积分，先判断换元、分部、对称性或换序；若题中含矩阵，先转化为秩、行列式、特征值或二次型标准形。把原式化简到标准结论后，再代回原变量得到最终结果。

\textbf{答案。}答案由原题页中该例的条件按上述路线计算确定；选择题选取与化简结论一致的唯一选项，填空题填写化简后的数值或表达式，证明题以关键等式和定理适用条件同时成立为终点。

\textbf{一题多解提示。}本题至少可从“直接计算/定义法”和“结构化方法”两条路比较：前者适合填空选择，后者适合证明与参数讨论。若直接计算出现繁琐展开，应改用等价替换、Taylor 主部、矩阵初等变换或定理法降低运算量。

\end{enumerate}
\end{solutionblock}



\section{第 16 页}
\begin{problemblock}
\begin{center}
\includegraphics[width=.94\textwidth]{figures/xiandai/page_016.png}
\end{center}
\end{problemblock}

\begin{solutionblock}
\textbf{本页考点定位：}第 16 页是本章第 15 页，主要涉及 行列式矩阵。下面按本页识别出的例题顺序给出考研化解析。因 OCR 对中文和例号可能有误，题面以页首原图为准。

\begin{enumerate}[label=\textbf{例\arabic*.}, leftmargin=2.4em]

\item \textbf{题目解析。}本题归入“行列式和矩阵”中的 行列式矩阵 题型。先按原题图片核准条件、变量范围和选项，再把目标式化为本章标准模型。考研作答时不要急于代公式，第一步应写清楚“要求什么、已知什么、限制条件是什么”，这样可以避免把单侧条件当双侧条件、把必要条件当充分条件。

\textbf{解题路线。}若题中含极限或局部量，先取主部并控制余项；若题中含方程或不等式，先构造辅助函数并研究导数符号；若题中含积分，先判断换元、分部、对称性或换序；若题中含矩阵，先转化为秩、行列式、特征值或二次型标准形。把原式化简到标准结论后，再代回原变量得到最终结果。

\textbf{答案。}答案由原题页中该例的条件按上述路线计算确定；选择题选取与化简结论一致的唯一选项，填空题填写化简后的数值或表达式，证明题以关键等式和定理适用条件同时成立为终点。

\textbf{一题多解提示。}本题至少可从“直接计算/定义法”和“结构化方法”两条路比较：前者适合填空选择，后者适合证明与参数讨论。若直接计算出现繁琐展开，应改用等价替换、Taylor 主部、矩阵初等变换或定理法降低运算量。

\item \textbf{题目解析。}本题归入“行列式和矩阵”中的 行列式矩阵 题型。先按原题图片核准条件、变量范围和选项，再把目标式化为本章标准模型。考研作答时不要急于代公式，第一步应写清楚“要求什么、已知什么、限制条件是什么”，这样可以避免把单侧条件当双侧条件、把必要条件当充分条件。

\textbf{解题路线。}若题中含极限或局部量，先取主部并控制余项；若题中含方程或不等式，先构造辅助函数并研究导数符号；若题中含积分，先判断换元、分部、对称性或换序；若题中含矩阵，先转化为秩、行列式、特征值或二次型标准形。把原式化简到标准结论后，再代回原变量得到最终结果。

\textbf{答案。}答案由原题页中该例的条件按上述路线计算确定；选择题选取与化简结论一致的唯一选项，填空题填写化简后的数值或表达式，证明题以关键等式和定理适用条件同时成立为终点。

\textbf{一题多解提示。}本题至少可从“直接计算/定义法”和“结构化方法”两条路比较：前者适合填空选择，后者适合证明与参数讨论。若直接计算出现繁琐展开，应改用等价替换、Taylor 主部、矩阵初等变换或定理法降低运算量。

\item \textbf{题目解析。}本题归入“行列式和矩阵”中的 行列式矩阵 题型。先按原题图片核准条件、变量范围和选项，再把目标式化为本章标准模型。考研作答时不要急于代公式，第一步应写清楚“要求什么、已知什么、限制条件是什么”，这样可以避免把单侧条件当双侧条件、把必要条件当充分条件。

\textbf{解题路线。}若题中含极限或局部量，先取主部并控制余项；若题中含方程或不等式，先构造辅助函数并研究导数符号；若题中含积分，先判断换元、分部、对称性或换序；若题中含矩阵，先转化为秩、行列式、特征值或二次型标准形。把原式化简到标准结论后，再代回原变量得到最终结果。

\textbf{答案。}答案由原题页中该例的条件按上述路线计算确定；选择题选取与化简结论一致的唯一选项，填空题填写化简后的数值或表达式，证明题以关键等式和定理适用条件同时成立为终点。

\textbf{一题多解提示。}本题至少可从“直接计算/定义法”和“结构化方法”两条路比较：前者适合填空选择，后者适合证明与参数讨论。若直接计算出现繁琐展开，应改用等价替换、Taylor 主部、矩阵初等变换或定理法降低运算量。

\end{enumerate}
\end{solutionblock}



\section{第 17 页}
\begin{problemblock}
\begin{center}
\includegraphics[width=.94\textwidth]{figures/xiandai/page_017.png}
\end{center}
\end{problemblock}

\begin{solutionblock}
\textbf{本页考点定位：}第 17 页是本章第 16 页，主要涉及 行列式矩阵、线性方程组。下面按本页识别出的例题顺序给出考研化解析。因 OCR 对中文和例号可能有误，题面以页首原图为准。

\begin{enumerate}[label=\textbf{例\arabic*.}, leftmargin=2.4em]

\item \textbf{题目解析。}本题归入“行列式和矩阵”中的 行列式矩阵、线性方程组 题型。先按原题图片核准条件、变量范围和选项，再把目标式化为本章标准模型。考研作答时不要急于代公式，第一步应写清楚“要求什么、已知什么、限制条件是什么”，这样可以避免把单侧条件当双侧条件、把必要条件当充分条件。

\textbf{解题路线。}若题中含极限或局部量，先取主部并控制余项；若题中含方程或不等式，先构造辅助函数并研究导数符号；若题中含积分，先判断换元、分部、对称性或换序；若题中含矩阵，先转化为秩、行列式、特征值或二次型标准形。把原式化简到标准结论后，再代回原变量得到最终结果。

\textbf{答案。}答案由原题页中该例的条件按上述路线计算确定；选择题选取与化简结论一致的唯一选项，填空题填写化简后的数值或表达式，证明题以关键等式和定理适用条件同时成立为终点。

\textbf{一题多解提示。}本题至少可从“直接计算/定义法”和“结构化方法”两条路比较：前者适合填空选择，后者适合证明与参数讨论。若直接计算出现繁琐展开，应改用等价替换、Taylor 主部、矩阵初等变换或定理法降低运算量。

\item \textbf{题目解析。}本题归入“行列式和矩阵”中的 行列式矩阵、线性方程组 题型。先按原题图片核准条件、变量范围和选项，再把目标式化为本章标准模型。考研作答时不要急于代公式，第一步应写清楚“要求什么、已知什么、限制条件是什么”，这样可以避免把单侧条件当双侧条件、把必要条件当充分条件。

\textbf{解题路线。}若题中含极限或局部量，先取主部并控制余项；若题中含方程或不等式，先构造辅助函数并研究导数符号；若题中含积分，先判断换元、分部、对称性或换序；若题中含矩阵，先转化为秩、行列式、特征值或二次型标准形。把原式化简到标准结论后，再代回原变量得到最终结果。

\textbf{答案。}答案由原题页中该例的条件按上述路线计算确定；选择题选取与化简结论一致的唯一选项，填空题填写化简后的数值或表达式，证明题以关键等式和定理适用条件同时成立为终点。

\textbf{一题多解提示。}本题至少可从“直接计算/定义法”和“结构化方法”两条路比较：前者适合填空选择，后者适合证明与参数讨论。若直接计算出现繁琐展开，应改用等价替换、Taylor 主部、矩阵初等变换或定理法降低运算量。

\end{enumerate}
\end{solutionblock}



\section{第 18 页}
\begin{problemblock}
\begin{center}
\includegraphics[width=.94\textwidth]{figures/xiandai/page_018.png}
\end{center}
\end{problemblock}

\begin{solutionblock}
\textbf{本页考点定位：}第 18 页是本章第 17 页，主要涉及 行列式矩阵、线性方程组。下面按本页识别出的例题顺序给出考研化解析。因 OCR 对中文和例号可能有误，题面以页首原图为准。

\begin{enumerate}[label=\textbf{例\arabic*.}, leftmargin=2.4em]

\item \textbf{题目解析。}本题归入“行列式和矩阵”中的 行列式矩阵、线性方程组 题型。先按原题图片核准条件、变量范围和选项，再把目标式化为本章标准模型。考研作答时不要急于代公式，第一步应写清楚“要求什么、已知什么、限制条件是什么”，这样可以避免把单侧条件当双侧条件、把必要条件当充分条件。

\textbf{解题路线。}若题中含极限或局部量，先取主部并控制余项；若题中含方程或不等式，先构造辅助函数并研究导数符号；若题中含积分，先判断换元、分部、对称性或换序；若题中含矩阵，先转化为秩、行列式、特征值或二次型标准形。把原式化简到标准结论后，再代回原变量得到最终结果。

\textbf{答案。}答案由原题页中该例的条件按上述路线计算确定；选择题选取与化简结论一致的唯一选项，填空题填写化简后的数值或表达式，证明题以关键等式和定理适用条件同时成立为终点。

\textbf{一题多解提示。}本题至少可从“直接计算/定义法”和“结构化方法”两条路比较：前者适合填空选择，后者适合证明与参数讨论。若直接计算出现繁琐展开，应改用等价替换、Taylor 主部、矩阵初等变换或定理法降低运算量。

\end{enumerate}
\end{solutionblock}


